\documentclass[11pt, a4paper]{article}

\usepackage[T1]{fontenc}
\usepackage[english]{babel}
\usepackage[
  top=3.2cm, bottom=3.2cm,
  left=2.8cm, right=2.8cm
]{geometry}
\usepackage{setspace}
\setstretch{1.15}
\setlength{\parskip}{0.5em}
\setlength{\parindent}{0pt}

\usepackage{xcolor}
\usepackage{hyperref}
\usepackage{booktabs}
\usepackage{enumitem}
\usepackage{listings}
\usepackage{amsmath}
\usepackage{fancyhdr}
\usepackage{titlesec}
\usepackage{caption}
\usepackage{float}
\usepackage{tabularx}

\definecolor{accent}{HTML}{2b5797}
\definecolor{linkblue}{HTML}{1a5276}
\definecolor{codebg}{HTML}{f4f4f8}
\definecolor{codegrey}{HTML}{555555}

\hypersetup{
  colorlinks=true,
  linkcolor=accent,
  citecolor=accent,
  urlcolor=linkblue,
  pdftitle={Tessera: The Permission Layer for AI Agents},
  pdfauthor={Guglielmo Reggio}
}

\titleformat{\section}{\Large\bfseries}{\thesection.}{0.6em}{}
\titleformat{\subsection}{\large\bfseries}{\thesubsection}{0.5em}{}
\titlespacing{\section}{0pt}{2em}{0.8em}
\titlespacing{\subsection}{0pt}{1.4em}{0.5em}

\lstset{
  backgroundcolor=\color{codebg},
  basicstyle=\small\ttfamily\color{codegrey},
  breaklines=true,
  frame=single,
  framerule=0pt,
  rulecolor=\color{codebg},
  xleftmargin=1em,
  xrightmargin=1em,
  aboveskip=1em,
  belowskip=1em,
  framexleftmargin=0.5em,
  framexrightmargin=0.5em,
  framextopmargin=0.5em,
  framexbottommargin=0.5em,
}

\pagestyle{fancy}
\fancyhf{}
\renewcommand{\headrulewidth}{0pt}
\fancyfoot[C]{\small\color{gray} Tessera --- Working Draft v0.6 --- April 2026 \hfill \thepage}

\setlist[itemize]{topsep=0.3em, itemsep=0.15em, parsep=0pt}
\setlist[enumerate]{topsep=0.3em, itemsep=0.15em, parsep=0pt}

\begin{document}

\thispagestyle{empty}
\vspace*{3.5cm}

\begin{center}
  {\Huge\bfseries TESSERA}\\[0.8em]
  {\Large The permission layer for AI agents}\\[2.5em]
  {\large Working Draft v0.6}\\[1.2em]
  {\large Guglielmo Reggio}\\[0.6em]
  {\normalsize April 2026}\\[4em]
  \rule{8cm}{0.4pt}\\[1.5em]
  {\itshape\small
    The Roman \emph{tessera hospitalis} was a token split between two strangers.\\
    Reuniting the pieces proved trust across distance.
  }
\end{center}

\newpage
\setcounter{page}{1}
\tableofcontents
\newpage

\begin{abstract}
\noindent AI agents can now send messages, execute code, spend money, publish content, and operate tools across the web. But the internet still lacks a standard way to answer four questions at the moment an agent acts: who authorised this agent, what is it allowed to do, when does that authority expire, and can that authority be revoked immediately?

Tessera is an open protocol for execution-time authorisation of AI agents. It makes sensitive agent actions contingent on a valid, human-delegated credential. Before an agent can send an email, trigger a payment, execute a shell command, or publish content, it must present cryptographic proof that a real human authorised that class of action within a defined scope and time window.

Tessera's first product is Tessera Guard: a runtime middleware layer that sits between agent frameworks and sensitive tools, verifies agent credentials, and blocks unauthorised actions by default.

Under the hood, Tessera combines a human root credential, privacy-preserving presentation, scoped delegation from human to agent, revocation and policy verification at execution time, and an open SDK for agent runtimes, tool gateways, and platforms.

Tessera is not a cryptocurrency, does not depend on biometric hardware, and does not require platforms to learn user identity. Its purpose is narrower and more practical: to become the standard authorisation layer for AI agents.
\end{abstract}

\section{The Problem}

\subsection{Agents now have real-world power}

AI systems are no longer limited to generating text. They increasingly act through agents that can:

\begin{itemize}
  \item send emails and messages,
  \item modify code and infrastructure,
  \item call APIs and external services,
  \item make purchases and subscriptions,
  \item post public content,
  \item and coordinate with other agents.
\end{itemize}

This is a structural shift. Once an agent can act, the key question is no longer just whether it is intelligent. The key question is whether it is authorised.

\subsection{Identity is not enough}

Traditional identity systems answer questions about who a user is. They do not cleanly answer what an autonomous agent is allowed to do on that user's behalf at the moment of execution.

Existing systems rely on one of three flawed patterns:

\begin{enumerate}
  \item \textbf{Ambient authority} --- if the agent has access to a tool or token, it can act.
  \item \textbf{Repeated approval prompts} --- the human is asked every time, destroying the utility of autonomy.
  \item \textbf{Platform-specific permissions} --- each runtime reinvents authorisation in its own format.
\end{enumerate}

None of these creates portable, auditable, revocable delegation.

\subsection{The missing primitive}

The internet needs a primitive for human-to-agent delegated authority.

That primitive must allow a verifier to check, in real time:

\begin{itemize}
  \item that a human principal exists,
  \item that the human delegated authority to this specific agent,
  \item that the delegated scope covers the requested action,
  \item that the authorisation is still valid,
  \item and that it has not been revoked.
\end{itemize}

That is the problem Tessera solves.

\section{What Tessera Is}

Tessera is an open authorisation protocol for agentic systems.

It introduces two credential types.

\subsection{Human credential}

A human credential proves that a real human principal exists and is eligible to delegate authority.

\subsection{Agent credential}

An agent credential is issued by that human to an agent, defining:

\begin{itemize}
  \item the agent identity,
  \item the permitted actions,
  \item optional spending or usage limits,
  \item optional resource constraints,
  \item the issuance time,
  \item the expiry time,
  \item and the human signature authorising the delegation.
\end{itemize}

The protocol is designed so that sensitive actions are checked at execution time, not assumed from prior login state or API possession.

\section{The First Product: Tessera Guard}

Tessera Guard is the first implementation of the protocol.

It is a middleware layer that intercepts sensitive tool calls and blocks them unless the agent presents a valid Tessera credential.

\subsection{Guarded action classes}

The initial protected classes are:

\begin{itemize}
  \item \texttt{message.send}
  \item \texttt{payment.intent}
  \item \texttt{exec.shell}
  \item \texttt{content.publish}
\end{itemize}

These classes are deliberately narrow. They cover high-risk actions with clear user harm and immediate practical demand.

\subsection{Verification flow}

When an agent attempts a guarded action:

\begin{enumerate}
  \item The runtime passes the requested action and attached credential to Tessera Guard.
  \item Guard verifies credential validity.
  \item Guard verifies issuer and delegation signatures.
  \item Guard checks expiry.
  \item Guard checks revocation status.
  \item Guard checks whether the requested action falls within delegated scope.
  \item Guard returns allow or deny, with a structured explanation.
\end{enumerate}

This turns authorisation into a deterministic runtime boundary.

\subsection{Why Guard matters}

Tessera Guard changes the default posture of agent systems:

\begin{itemize}
  \item without Guard: agents act unless stopped
  \item with Guard: agents cannot act unless authorised
\end{itemize}

That inversion is the core product value.

\section{Design Principles}

Tessera is designed around the following principles.

\subsection{Execution-time authority}

Authority must be evaluated when the action executes, not inferred from long-lived access.

\subsection{Scoped delegation}

An agent can act only within the authority explicitly granted to it.

\subsection{Fast revocation}

A human must be able to invalidate delegated authority immediately.

\subsection{Minimal disclosure}

Verifiers should learn only what they need to know to authorise the action.

\subsection{Runtime portability}

Authorisation should work across agent frameworks, MCP servers, gateways, and platforms.

\subsection{Open protocol, not closed platform lock-in}

The core credential and verification model should be implementable by anyone.

\section{Protocol Model}

Tessera has four layers.

\subsection{Layer 1 --- Human root credential}

A human obtains a root credential from an issuer. This establishes eligibility to delegate authority.

The exact trust anchor may vary by deployment and jurisdiction. In early versions, Tessera may rely on existing verification rails such as regulated identity or financial verification providers. The protocol does not require revealing personal identity to verifiers during runtime authorisation.

\subsection{Layer 2 --- Agent identity}

Each agent has its own public key and stable identifier.

An agent is not merely a session. It is a cryptographic subject capable of receiving delegated authority.

\subsection{Layer 3 --- Delegation}

The human signs a delegation to the agent. The delegation includes:

\begin{itemize}
  \item agent public key or identifier,
  \item allowed action classes,
  \item optional quantitative limits,
  \item optional contextual constraints,
  \item issuance time,
  \item expiry time,
  \item revocation handle,
  \item human signature.
\end{itemize}

\subsection{Layer 4 --- Verification}

At runtime, a verifier checks:

\begin{itemize}
  \item credential integrity,
  \item signature chain,
  \item scope match,
  \item expiration,
  \item revocation,
  \item and policy constraints.
\end{itemize}

The result is a simple authorisation decision.

\section{Credential Structure}

\subsection{Human credential}

A Tessera human credential attests that the holder is authorised to delegate agent authority.

The verifier should not need to know the person's name, bank, email, or legal identity to evaluate a delegated action. The runtime only needs to know whether the human root is valid and what assurance tier or policy class it carries.

\subsection{Agent credential}

A Tessera agent credential contains:

\begin{itemize}
  \item a reference to the parent human root,
  \item the agent identifier,
  \item the delegation scope,
  \item issuance and expiry,
  \item and the parent signature.
\end{itemize}

Example shape:

\begin{lstlisting}
{
  "type": "TesseraAgentCredential",
  "agent": "agent_pk_123",
  "parent": "human_commitment_abc",
  "scope": {
    "actions": ["message.send", "content.publish"],
    "max_recipients": 20,
    "allowed_domains": ["example.com"]
  },
  "issued_at": 1743465600,
  "expires_at": 1743552000,
  "revocation_id": "rev_456",
  "signature": "..."
}
\end{lstlisting}

\subsection{Containment rule}

No child delegation may exceed the permissions of its parent.

This allows safe sub-delegation in multi-agent systems while preserving a clear authority chain.

\section{Revocation}

Revocation is a first-class requirement.

If a human revokes an agent credential, the next guarded action must fail immediately.

Tessera therefore treats revocation not as a slow administrative process, but as an online part of authorisation.

A production deployment may use a combination of:

\begin{itemize}
  \item status registries,
  \item signed status lists,
  \item cached freshness windows,
  \item and high-availability revocation endpoints.
\end{itemize}

The exact implementation can evolve, but the product requirement is fixed:

\begin{quote}
revocation must be fast enough to function as a real kill switch.
\end{quote}

\section{Privacy Model}

Tessera aims to minimise disclosure during verification.

A verifier should learn:

\begin{itemize}
  \item whether the credential is valid,
  \item whether the action is within scope,
  \item when the credential expires,
  \item and what policy tier applies.
\end{itemize}

A verifier should not need to learn:

\begin{itemize}
  \item the human's civil identity,
  \item the underlying verification source,
  \item unrelated activity history,
  \item or linkable information across unrelated verifications.
\end{itemize}

Privacy is not the entire product, but it matters. Authorisation for agents should not require turning every delegated action into a full identity disclosure event.

\section{Threat Model}

Tessera is designed to reduce the following risks:

\begin{table}[H]
\centering
\small
\begin{tabularx}{\textwidth}{@{}lX@{}}
\toprule
\textbf{Risk} & \textbf{Defence} \\
\midrule
Unauthorised action & Execution-time scope checks. \\
Credential replay & Bounded validity, agent-bound keys, and revocation. \\
Scope escalation & Deterministic scope evaluation and containment rules. \\
Multi-agent ambiguity & Signed delegation chain from human root to acting agent. \\
Opaque failure & Structured denial reasons that agents can surface to users. \\
\bottomrule
\end{tabularx}
\caption{Core Tessera threat model.}
\end{table}

Tessera does not solve every misuse problem. A real human may still authorise harmful behaviour. Tessera's role is narrower: it makes authority explicit, bounded, and auditable.

\section{Integration Surfaces}

Tessera is designed for four initial integration surfaces.

\subsection{Agent runtimes}

Open-source and commercial runtimes can add Guard to their tool invocation path.

\subsection{MCP servers and gateways}

Servers that expose tools to agents can require Tessera credentials before serving sensitive operations.

\subsection{Developer platforms and APIs}

Services can distinguish between unauthenticated automation and human-authorised agent actions.

\subsection{Enterprise internal agents}

Organisations can require bounded delegation for agents acting in internal systems.

\section{Why Now}

Three trends make this the right moment.

\subsection{Agents are crossing from assistance into action}

The risk surface has moved from text generation to real-world execution.

\subsection{Security models are lagging}

Most agent systems still rely on ambient tool access, brittle prompts, or repeated approval loops.

\subsection{Standards are still open}

The market is early enough that the authorisation layer for agents has not yet ossified into a single closed vendor model.

This creates room for an open protocol to define the category.

\section{Go-to-Market}

Tessera starts where the pain is highest and the feedback loop is fastest.

\subsection{Phase 1 --- OpenClaw and open agent ecosystems}

The first battlefield is the open agent developer community.

Why:

\begin{itemize}
  \item these users already understand tools and permissions,
  \item they can adopt middleware quickly,
  \item they produce visible design feedback,
  \item and they are highly valuable for protocol iteration.
\end{itemize}

The goal of this phase is not maximum revenue. It is product truth.

\subsection{Phase 2 --- Cross-runtime expansion}

Once validated in one ecosystem, Tessera expands to:

\begin{itemize}
  \item MCP-compatible tools,
  \item agent gateways,
  \item coding agents,
  \item workflow automations,
  \item and API services receiving agent traffic.
\end{itemize}

\subsection{Phase 3 --- Enterprise control plane}

The long-term commercial layer includes:

\begin{itemize}
  \item hosted verification,
  \item revocation infrastructure,
  \item policy administration,
  \item audit trails,
  \item team and org controls,
  \item and enterprise integrations.
\end{itemize}

This creates a natural model:

\begin{itemize}
  \item open protocol,
  \item open-source Guard,
  \item managed control plane.
\end{itemize}

\section{Business Model}

Tessera's commercial model should remain tightly aligned to authorisation infrastructure.

The clearest monetisable services are:

\begin{itemize}
  \item hosted verification,
  \item revocation infrastructure,
  \item policy management,
  \item enterprise auditability,
  \item and premium support.
\end{itemize}

This keeps the business legible. The protocol can remain open while the operational control plane becomes the commercial offering.

Early versions should resist adding unnecessary monetary complexity. The main value is not speculative economics. The main value is safe delegated execution.

\section{Non-goals for the Initial Version}

Tessera is not, in its initial form:

\begin{itemize}
  \item a universal social network identity,
  \item a cryptocurrency,
  \item a consumer payments system,
  \item a full legal identity replacement,
  \item or a generalised provenance system for all digital content.
\end{itemize}

Those may become adjacent extensions over time, but they are not the launch product.

The launch product is narrower:

\begin{quote}
authorisation for sensitive AI agent actions.
\end{quote}

\section{Roadmap}

\begin{table}[H]
\centering
\small
\begin{tabularx}{\textwidth}{@{}lX@{}}
\toprule
\textbf{Phase} & \textbf{Scope} \\
\midrule
Phase 0 --- Protocol and SDK & Credential schema. Signature and verification logic. Local demo verifier. Basic revocation model. \\
Phase 1 --- Tessera Guard for OpenClaw & Tool interception middleware. Initial guarded action classes. Structured denial responses. Real developer feedback. \\
Phase 2 --- Cross-runtime support & MCP and gateway integrations. Policy templates. Hosted verification service. \\
Phase 3 --- Enterprise layer & Multi-user admin. Team delegation. Audit and compliance features. High-availability revocation and policy infrastructure. \\
Phase 4 --- Broader standards participation & Standards engagement. Interop work. Multi-issuer ecosystem. \\
\bottomrule
\end{tabularx}
\caption{Narrowed v0.6 roadmap.}
\end{table}

\section{Why the Name Tessera}

A Roman \emph{tessera hospitalis} was a token broken into matching parts between strangers. Reuniting the pieces established trust across distance.

That image captures the protocol well: human authority on one side, agent action on the other, joined only when the pieces match.

\section{Conclusion}

AI agents need more than prompts, approvals, or platform-specific permissions. They need a standard way to prove who authorised them, what they are allowed to do, and whether that authority still holds.

Tessera exists to provide that layer.

Its purpose is simple:

\begin{quote}
before an agent can act, it must be authorised.
\end{quote}

That should become a standard property of the internet.

\section{Get Involved}

Tessera is looking for:

\begin{itemize}
  \item early runtime integrations,
  \item security and identity reviewers,
  \item design partners in agent infrastructure,
  \item and contributors to the open SDK and Guard.
\end{itemize}

The protocol is open because agent authorisation should become infrastructure, not a silo.

\vspace{1.5em}
\begin{center}
  GitHub: \url{https://github.com/tessera-protocol/tessera}
\end{center}

\vfill
\begin{center}
  \small\color{gray}
  Tessera Working Draft v0.6 --- April 2026 --- Licensed under CC BY 4.0
\end{center}

\end{document}
